\documentclass[a4paper,11pt,fleqn]{article}%fleqn pour aligner les equations à gauche
\usepackage{../../Styles/Cours}


\newcommand{\chnum}{Chapitre 1}
\newcommand{\chtitle}{Composition initiale d'un système}
\newcommand{\dtype}{Fiche Révision}

\begin{document}
	\maketitle
	
	\section*{Ce qu'il faut connaitre :}
	\begin{multicols}{2}
		\subsection*{Les grandeurs :}
		\begin{itemize}[label=$\square$]
			\item Constante d'Avogadro
			\item Masse molaire
			\item Volume molaire
			\item Concentration (en quantité de matière)
		\end{itemize}
		\subsection*{Les relations :}
		\begin{itemize}[label=$\square$]
			\item Relation entre masse molaire d’une espèce, masse des entités et constante d’Avogadro.
			\item Masse molaire atomique d’un élément.
			\item  Volume molaire d’un gaz.
			\item Concentration en quantité de matière.
		\end{itemize}
	\end{multicols}
	
	\section*{Ce qu'il faut savoir faire :}
	\begin{table}[h!]
		\begin{tabular}{|m{.6\linewidth}|m{.2\linewidth}|m{.1\linewidth}|}
			\hline
			\textbf{La compétence} & \textbf{Où la trouver} & \textbf{Je la maîtrise ?} \\
			\hline
			Déterminer la masse molaire d’une espèce à partir des masses molaires atomiques des éléments qui la composent. & & \\
			\hline
			Déterminer la quantité de matière contenue dans un échantillon de corps pur à partir de sa masse et du tableau périodique.& & \\
			\hline
			Utiliser le volume molaire d’un gaz pour déterminer une quantité de matière. Déterminer la quantité de matière de chaque espèce dans un mélange (solide ou liquide) à partir de sa composition.
			& & \\
			\hline
			Déterminer la quantité de matière d’un soluté à partir de sa concentration en masse ou en quantité de matière et du volume de solution. & & \\
			\hline
			%Expliquer ou prévoir la couleur d’une espèce en solution à partir de son spectre UV-Visible.
			
			
			%Déterminer la concentration d’un soluté à partir de données expérimentales relatives à l’absorbance de solutions de concentrations connues.
			%
			%Proposer et mettre en œuvre un protocole pour réaliser une gamme étalon et déterminer la concentration d’une espèce colorée en solution par des mesures d’absorbance. Tester les limites d’utilisation du protocole.& & \\
			%\hline
			
		\end{tabular}
	\end{table}
	
\end{document}
